\documentclass{article}
\newcommand{\shade}[1]{#1}
\newtheorem{theorem}{Theorem}
\begin{document}
A point is \shade{balanced} exactly when its two coordinate sums agree.
The ambient box is $B=[-2,2]^2$.
A later paragraph mentions balance in an unrelated accounting example.
\begin{theorem}\label{thm:balanced}
Every balanced point in the ambient box lies on the comparison locus.
\end{theorem}
\begin{proof}
The equality of the coordinate sums gives the comparison-locus equation inside $B$.
\end{proof}
\end{document}
